\documentclass[12pt]{article}
\usepackage{amsmath, amssymb, geometry, enumitem}
\geometry{margin=1in}
\setlength{\parskip}{0.6em}
\setlength{\parindent}{0pt}

\begin{document}

\begin{center}
\Large \textbf{ECON 300 -- Labour Economics} \vspace{.5cm}\\
\large \textbf{Quiz 5} \vspace{.5cm}\\
\large \textbf{April 07, 2026}
\end{center}

\begin{enumerate}



\item A country admits \textbf{420,000} immigrants in one year. They are divided into three categories:

\begin{itemize}
    \item Economic immigrants: \textbf{189,000}
    \item Family-class immigrants: \textbf{147,000}
    \item Refugees and humanitarian immigrants: \textbf{84,000}
\end{itemize}


\begin{enumerate}[label=(\alph*)]
    \item Calculate the percentage share of each category in total immigration. \hfill 

    \item Next year, the government increases economic immigrants by \textbf{18\%}, reduces family-class immigrants by \textbf{8\%}, and leaves refugees/humanitarian immigrants unchanged. Calculate the new number in each category and the new total number of immigrants. \hfill 

    \item Under the new totals, calculate the new percentage share of economic immigrants. \hfill 
\end{enumerate}



\item In a labour market, the working-age population is \textbf{6,800,000}. Initially:

\begin{itemize}
    \item Labour force participation rate = \textbf{67\%}
    \item Unemployment rate = \textbf{6.5\%}
\end{itemize}



\begin{enumerate}[label=(\alph*)]
    \item Calculate the size of the labour force. \hfill 

    \item Calculate the number of unemployed workers. \hfill 

    \item Calculate the number of employed workers. \hfill 

    \item Calculate the employment rate. \hfill 
\end{enumerate}

Now suppose that during a recession:

\begin{itemize}
    \item \textbf{52,000} employed workers lose their jobs and begin searching
    \item \textbf{36,000} unemployed workers become discouraged and leave the labour force
    \item \textbf{18,000} people previously out of the labour force enter the labour force and immediately begin searching for work
\end{itemize}

\begin{enumerate}[label=(\alph*), start=5]
    \item Calculate the new number of unemployed workers. \hfill 

    \item Calculate the new labour force. \hfill 

    \item Calculate the new unemployment rate. Briefly state whether the unemployment rate fully reflects the deterioration in labour market conditions. \hfill 
\end{enumerate}


\item Two countries, A and B, each have a labour force of \textbf{12,500,000}.

\begin{itemize}
    \item Country A has an unemployment rate of \textbf{5.6\%}
    \item Country B has an unemployment rate of \textbf{5.6\%}
\end{itemize}

However, the composition of unemployment differs:

\begin{itemize}
    \item In Country A, \textbf{14\%} of the unemployed have been unemployed for \textbf{12 months or more}
    \item In Country B, \textbf{41\%} of the unemployed have been unemployed for \textbf{12 months or more}
\end{itemize}



\begin{enumerate}[label=(\alph*)]
    \item Calculate the number of unemployed workers in each country. \hfill 

    \item Calculate the number of long-term unemployed workers in each country. \hfill 

    \item Calculate the long-term unemployment rate in each country as a percentage of the labour force. \hfill 

    \item Both countries have the same overall unemployment rate. Explain briefly why Country B may still have a more serious labour market problem. \hfill 
\end{enumerate}
\end{enumerate}
\end{document}